\clearpage
\pagenumbering{roman}
\chapter*{Preface}
\addcontentsline{toc}{chapter}{Preface}

This book is written in English. It follows two intellectual threads---mathematics and artificial intelligence---along one narrative spine:

\[
\text{Representation}\ \rightarrow\ \text{Learning}\ \rightarrow\ \text{Generation}\ \rightarrow\ \text{Intelligence}.
\]

Mathematics is the language. AI is the primary laboratory. Other sciences---finance, physics, control, information theory---enter only when a shared structure is genuine.

The project around this book also ships a website (PDF reader), Manim animations, slide decks, and small examples. Those artifacts are not substitutes for the text; they illustrate it.

\chapter*{How to Read}
\addcontentsline{toc}{chapter}{How to Read}

Each chapter asks three questions:

\begin{enumerate}
  \item What is the underlying mathematical structure?
  \item Why does that structure appear in this model or algorithm?
  \item Where else does the same structure appear?
\end{enumerate}

Draft chapters mark what is still open. Speculation is never presented as fact.
